\documentclass[11pt]{article}
\usepackage[a4paper,margin=1in]{geometry}
\usepackage{subfiles}
\usepackage{amsmath,amssymb,bm}
\usepackage{hyperref}
\usepackage{booktabs}
\usepackage{array}
\usepackage{mathtools}

\title{Appendix 3 (to main theory): Emergent Bi--Conformal Geometry from the ISPG Pressure Ontology}
\author{}
\date{}

\begin{document}
\let\phi\varphi
\maketitle

\section*{Purpose and scope}
This appendix states the \emph{logic} that leads from the ISPG ontology (space as a physical medium with a static pressure field) to the \emph{bi--conformal} spacetime geometry used throughout the main theory. The goal is to make the emergence of the metric transparent, to clarify the one--metric (minimal coupling) posture, and to record the minimal assumptions under which the bi--conformal form is uniquely selected.

\paragraph{Non-negotiable posture.}
The main theory is formulated with a \emph{single} physical metric $g_{\mu\nu}$ to which matter is minimally coupled. No Einstein/Jordan frame change and no disformal matter metric is introduced. The ``emergent geometry'' content is therefore the statement that, in the ISPG programme, $g_{\mu\nu}$ is not an arbitrary 10-component field but is \emph{reduced} to a one-scalar family $g_{\mu\nu}[\phi]$ determined by the pressure potential $\phi$.

\section{Pressure potential and operational time/length standards}\label{sec:operational}
\subsection{Pressure potential}
The central scalar is the dimensionless pressure potential
\begin{equation}
\phi \;\equiv\; \ln\!\left(\frac{P_{\rm stat}}{P_{\max}}\right),
\label{eq:phidef}
\end{equation}
so that pressure ratios add under composition and $\phi=0$ is the asymptotic vacuum normalization $P_{\rm stat}=P_{\max}$.

\subsection{Operational postulates (clocks and rods)}
ISPG fixes geometry by specifying how \emph{operational standards} respond to the medium state. The emergence of the metric rests on three minimal postulates:

\begin{enumerate}
\item \textbf{Clock-rate postulate (time as process tempo).}
The squared proper-time interval of a local resonant matter process is proportional to the local static pressure,
\begin{equation}
d\tau^2 \;\propto\; P_{\rm stat}\,dt^2
\quad\Longrightarrow\quad
d\tau \;=\; e^{\phi/2}\,dt ,
\label{eq:clock}
\end{equation}
after choosing the normalization $\tau=t$ at $\phi=0$. The squared form is the natural metric-level statement ($|g_{tt}|\propto P_{\rm stat}$), and the $1/2$ exponent follows directly by taking the square root. This is an operational law for clocks, atoms, and localized resonant processes; the medium itself is not assigned a separate time variable. Slowing is only an operational property of resonant clocks and processes. The same exponent is fixed independently in the main theory (\S\ref{sec:c_invariance}) by the requirement that the locally measured speed of light remains invariant ($c_{\rm meas}=c$) while the coordinate speed responds to the medium state; the two routes coincide and select the clock law uniquely.

\item \textbf{Reciprocal scale postulate (rod contraction).}
The same pressure state also rescales spatial rulers reciprocally so that local physics remains self-consistent under uniform changes of $P_{\rm stat}$. The rod law is a co-response to the pressure scalar, not a temporal slowing of space itself. In particular, reciprocal rescaling is the unique choice ensuring that locally measured velocities (including the speed of light) remain invariant under changes of the medium state, as derived in the main theory (\S\ref{sec:c_invariance}). This is implemented as an isotropic spatial scale factor
\begin{equation}
d\ell = e^{-\phi/2}\,d\sigma,
\label{eq:rod}
\end{equation}
where $d\sigma^2$ is the Euclidean line element in the chosen coordinates.

\item \textbf{Isotropy and locality.}
In the static exterior of an isolated source the medium state is spherically symmetric and depends only on $r$, so the rescaling must be isotropic in the spatial directions.
\end{enumerate}

Together, \eqref{eq:clock}--\eqref{eq:rod} define a unique (up to coordinate relabeling) static, isotropic metric.

\paragraph{Active-existence reading.}
The same postulates also admit a useful operational reading.  Local clocks
and rods compare relations among interaction-active excitation nodes of the
same medium; they do not directly measure an independently observable empty
interval between such nodes.  In later mass-sector notation this common
pressure response may be denoted by an operational ratio
$\mathcal N_{\rm act}(\phi)/\mathcal N_{\rm act}(0)=e^{\phi/2}$.
Here $\mathcal N_{\rm act}$ is not an independent number-density count or a
three-dimensional count of medium nodes; it is only a shorthand for the
already stated pressure-scaling ratio of matter standards.  This
interpretation does not add a third operational postulate and does not replace
the clock/rod proof of local light-speed invariance.

\section{Derivation of the bi--conformal line element}\label{sec:derive}
Combining \eqref{eq:clock} and \eqref{eq:rod} gives
\begin{equation}
ds^2 \;\equiv\; -c^2d\tau^2 + d\ell^2
\;=\; -e^{\phi}c^2dt^2 + e^{-\phi}d\sigma^2,
\label{eq:biconf_app3}
\end{equation}
with
\begin{equation}
d\sigma^2 \equiv dx^2+dy^2+dz^2
\quad\text{(or in spherical symmetry }d\sigma^2=dr^2+r^2d\Omega^2\text{)}.
\end{equation}
Equation~\eqref{eq:biconf_app3} is the \textbf{bi--conformal metric} used in the main theory. It is ``bi--conformal'' because the time and space parts scale with opposite exponents of the same scalar $\phi$.

\paragraph{Uniqueness at this level.}
Given (i) a single scalar medium state $\phi$, (ii) isotropy, and (iii) reciprocal rescaling of clocks and rods, \eqref{eq:biconf_app3} is the unique static diagonal metric with one free function $\phi(r)$.

\section{Stationary rotating extension}\label{sec:stationary}
Equation~\eqref{eq:biconf_app3} is the non-rotating limit of the one-metric geometry. For stationary rotating configurations one must allow an off-diagonal gravitomagnetic sector $g_{0i}$ while keeping the same scalar-controlled diagonal factors. To leading order in slow rotation, the natural stationary extension is
\begin{equation}
ds^2 \;=\; -e^{\phi}c^2dt^2 \;-\; \frac{4}{c^2}\,(\bm{A}\!\cdot\! d\bm{x})\,dt \;+\; e^{-\phi}d\sigma^2 + O(A^2),
\label{eq:biconf_stationary}
\end{equation}
where $\bm{A}$ is the gravitomagnetic vector potential.

\paragraph{Interpretational consequence.}
The scalar $\phi$ continues to encode the static pressure state (clock rate, rod scale, Newtonian potential), while rotational physics lives in the off-diagonal $g_{0i}$ sector. Therefore any MOND-motivated enhancement tied specifically to space rotation should enter through the constitutive law of $\bm{A}$ or $g_{0i}$, not by altering the diagonal bi-conformal core \eqref{eq:biconf_app3}.

\section{Weak-field limit and Newtonian correspondence}\label{sec:weak}
For small $|\phi|\ll 1$,
\begin{equation}
g_{tt}=-e^{\phi}\simeq -(1+\phi),\qquad
g_{ij}=e^{-\phi}\delta_{ij}\simeq (1-\phi)\delta_{ij}.
\end{equation}
Comparing to standard weak-field notation
\begin{equation}
g_{tt}\simeq -(1+2\Phi_N/c^2),\qquad g_{ij}\simeq (1-2\Psi/c^2)\delta_{ij},
\end{equation}
the bi--conformal form enforces
\begin{equation}
\Phi_N=\Psi=\frac{c^2}{2}\phi.
\label{eq:PhiPsi}
\end{equation}
Thus the metric automatically gives $\gamma_{\rm PPN}=\Psi/\Phi_N=1$ at leading order (a cornerstone consistency condition used throughout the main theory).

\section{One-metric minimal coupling and why no ``frame'' language is needed}\label{sec:onemetric}
In the main theory, matter is minimally coupled to $g_{\mu\nu}$ of \eqref{eq:biconf_app3}. The bi--conformal metric is \emph{not} a conformal redefinition of a separate ``Einstein-frame'' metric; it is the physical metric defined by the operational postulates. Consequently:
\begin{itemize}
\item there is no independent matter metric $\tilde g_{\mu\nu}$ in the main theory package;
\item arguments about propagation speeds and stability are made directly in the one-metric formulation (e.g.\ $\alpha_T=0$ for GW170817 consistency, principal-symbol hyperbolicity for causal propagation).
\end{itemize}

\section{Consistency with the covariant backbone and the sign issue}\label{sec:sign}
The main theory employs the scalar--tensor action
\begin{equation}
S=\frac{1}{16\pi G}\int d^4x\sqrt{-g}\left[R+\frac12 g^{\mu\nu}\partial_\mu\phi\,\partial_\nu\phi\right]+S_m[g_{\mu\nu},\psi],
\label{eq:action_app3}
\end{equation}
with the same single metric $g_{\mu\nu}$ as in \eqref{eq:biconf_app3}. The bi--conformal construction ties the scalar sector to the geometry sufficiently strongly that the field equations take the specific coupling form used in the main theory (including the $\kappa=-1/2$ factor in $G_{\mu\nu}=\kappa\,T^{(\phi)}_{\mu\nu}$, as verified componentwise in the main technical appendices). Appendix~11 then shows that the would-be ghost associated with the sign choice is eliminated by the constraint structure of the bi--conformal construction.

\paragraph{Interpretation.}
The ISPG posture is that \eqref{eq:biconf_app3} is the \emph{geometric encoding} of the medium state. The covariant action \eqref{eq:action_app3} is an efficient (EFT-level) way to compute the consequences of that geometry and to interface with standard tests (PPN/2PN, lensing, GW sector), while maintaining one-metric minimal coupling.

\section{Derived vs.\ stated: compact ledger}\label{sec:ledger_app3}
\begin{center}
\begin{tabular}{@{}p{0.28\linewidth}p{0.62\linewidth}@{}}
\toprule
\textbf{Derived here} &
Bi--conformal metric \eqref{eq:biconf_app3} from operational postulates \eqref{eq:clock}--\eqref{eq:rod}; weak-field identification \eqref{eq:PhiPsi} implying $\gamma_{\rm PPN}=1$ at leading order; one-metric (minimal coupling) posture.\\
\midrule
\textbf{Used in the main theory} &
Covariant backbone \eqref{eq:action_app3} as the EFT interface; post-Newtonian and wave-sector analyses built on \eqref{eq:biconf_app3}; constraint-based ghost elimination (Appendix~11).\\
\bottomrule
\end{tabular}
\end{center}

\section*{Takeaway}
ISPG fixes spacetime geometry by specifying how physical clocks and rods respond to the medium's static pressure state. With the pressure potential $\phi=\ln(P_{\rm stat}/P_{\max})$, the minimal operational postulates lead uniquely to the bi--conformal line element
\[
ds^2=-e^{\phi}c^2dt^2+e^{-\phi}d\sigma^2,
\]
which enforces $\Phi=\Psi$ in the weak field and underpins the one-metric, minimal-coupling formulation of the main theory. For rotating stationary systems, the same one-metric posture extends by adding an off-diagonal gravitomagnetic sector $g_{0i}$, while leaving the diagonal bi-conformal structure intact.

\end{document}
